\documentclass[11pt,a4paper]{article}

% ---------------------------------------------------------
% PREAMBLE (stable, warning-free, Overleaf-safe)
% ---------------------------------------------------------
\usepackage[utf8]{inputenc}
\usepackage[T1]{fontenc}
\usepackage{lmodern}
\usepackage{amsmath, amssymb, amsfonts}
\usepackage{bm}
\usepackage{geometry}
\usepackage{graphicx}
\usepackage{hyperref}
\usepackage{cite}
\usepackage{authblk}
\usepackage{physics}
\usepackage{mathtools}

\geometry{margin=2.4cm}

\title{\textbf{Companion XXXVIII: The Circumgalactic Medium in the Relaxation-Driven Cyclic Cosmology}}
\author{Michael Lehmann \\ \texttt{mi.lehmann@gmx.de}}
\date{April 2026}

\begin{document}
\maketitle

\begin{abstract}
The circumgalactic medium (CGM) forms the interface between galaxies and the cosmic 
web. It regulates gas inflows, outflows, and the long-term evolution of galaxies. 
In the Relaxation-Driven Cyclic Cosmology (RDCC), the scale-dependent evolution of 
the gravitational potential modifies the thermodynamic structure, kinematics, and 
ionization state of the CGM. This Companion develops a continuous description of 
the RDCC circumgalactic medium, deriving how the relaxation scale influences the 
density and temperature profiles, the stability of gaseous halos, and the dynamics 
of inflows and outflows. The resulting framework provides a natural explanation for 
the observed coherence of CGM absorption systems and the extended gaseous envelopes 
around massive galaxies.
\end{abstract}
% ---------------------------------------------------------
\section{Introduction}

The circumgalactic medium is a vast reservoir of gas that surrounds galaxies and 
mediates the exchange of mass, energy, and metals between galaxies and their 
environment. Observations reveal a multiphase structure extending hundreds of 
kiloparsecs, with cold filaments, warm ionized gas, and hot virialized atmospheres 
coexisting in a delicate balance.

In the standard cosmological model, the CGM is shaped by the interplay between 
gravitational collapse, radiative cooling, and feedback. In RDCC, the gravitational 
potential evolves differently. The relaxation mechanism suppresses the growth of 
small-scale modes and enhances the coherence of large-scale modes. This modifies 
the density and temperature structure of the CGM, the stability of gaseous halos, 
and the dynamics of inflows and outflows. The CGM becomes a direct tracer of the 
relaxation scale $k_{\mathrm{rel}}$.
% ---------------------------------------------------------
\section{The RDCC Potential and the Structure of Gaseous Halos}

The gravitational potential that shapes the CGM is modified in RDCC according to
\begin{equation}
    \Phi_{\mathrm{RDCC}}(k,a)
    = \Phi(k,a)
      \left[1 - \chi_{\Phi}
      \left(\frac{k}{k_{\mathrm{rel}}}\right)^{\alpha}\right].
\end{equation}

This modification reduces the depth of the potential wells of low-mass halos while 
preserving the coherence of the potential in massive halos. The hydrostatic 
equilibrium condition,
\begin{equation}
    \frac{1}{\rho}\frac{dP}{dr}
    = -\frac{d\Phi_{\mathrm{RDCC}}}{dr},
\end{equation}
then yields a density profile that is smoother and more extended than in the 
standard cosmological model.

The temperature profile follows from the ideal-gas relation,
\begin{equation}
    T_{\mathrm{RDCC}}(r)
    = \frac{\mu m_{p}}{k_{B}}
      \frac{1}{\rho(r)}
      \int_{r}^{\infty}
      \rho(r')\,\frac{d\Phi_{\mathrm{RDCC}}}{dr'}\,dr'.
\end{equation}

The resulting CGM is less centrally concentrated and exhibits a broader transition 
between the inner halo and the outer gaseous envelope.
% ---------------------------------------------------------
\section{Cooling, Condensation, and the Multiphase Structure}

The CGM is inherently multiphase, with cold clouds embedded in a hot background. 
The formation of these cold clouds depends on the competition between radiative 
cooling and gravitational compression. In RDCC, the cooling time becomes
\begin{equation}
    t_{\mathrm{cool,RDCC}}
    = \frac{3 k_{B} T_{\mathrm{RDCC}}}{n_{e} \Lambda(T_{\mathrm{RDCC}})},
\end{equation}
where the modified temperature profile delays cooling in low-mass halos and 
accelerates it in massive halos.

The condition for thermal instability,
\begin{equation}
    \frac{d\ln \Lambda}{d\ln T} < 2,
\end{equation}
is reached at different radii in RDCC than in the standard cosmological model. 
Massive halos develop extended regions of thermal instability, leading to the 
formation of cold filaments that penetrate deep into the halo. Low-mass halos, 
whose potentials are weakened by the relaxation mechanism, suppress the formation 
of cold clouds altogether.

This asymmetry provides a natural explanation for the observed scarcity of cold 
gas around dwarf galaxies and the prevalence of cold filaments around massive 
galaxies.
% ---------------------------------------------------------
\section{Inflow and Outflow Dynamics}

Gas flows into and out of galaxies through the CGM. In RDCC, the dynamics of these 
flows are modified by the relaxation-modified potential. The inflow velocity can be 
approximated as
\begin{equation}
    v_{\mathrm{in,RDCC}}(r)
    = \sqrt{
        \frac{2 G M_{\mathrm{RDCC}}(r)}{r}
    },
\end{equation}
where the reduced mass profile of low-mass halos slows inflows and the enhanced 
coherence of massive halos accelerates them.

Outflows driven by stellar or AGN feedback must overcome the escape velocity,
\begin{equation}
    v_{\mathrm{esc,RDCC}}(r)
    = \sqrt{
        2 \abs{\Phi_{\mathrm{RDCC}}(r)}
    }.
\end{equation}

Because the potential is shallower in low-mass halos, outflows escape more easily, 
leading to a CGM that is metal-poor and diffuse. Massive halos retain their 
outflows, enriching their CGM and promoting the formation of multiphase structures.

The CGM thus becomes a dynamical record of the relaxation scale.
% ---------------------------------------------------------
\section{Ionization Structure and Observational Signatures}

The ionization state of the CGM is sensitive to its density and temperature 
structure. In RDCC, the modified profiles lead to characteristic absorption 
signatures. The column density of an ion species $X$ becomes
\begin{equation}
    N_{X,\mathrm{RDCC}}
    = \int n_{X}(r)\,dr,
\end{equation}
where $n_{X}(r)$ reflects the RDCC-modified density and temperature.

Massive halos exhibit extended absorption systems with coherent velocity 
structures, while low-mass halos show weaker and more fragmented absorption. 
These predictions align with observations of Lyman-$\alpha$ absorbers and metal 
lines around galaxies across cosmic time.
% ---------------------------------------------------------
\section{Conclusion}

The circumgalactic medium in the Relaxation-Driven Cyclic Cosmology reflects the 
scale-dependent evolution of the gravitational potential. The relaxation mechanism 
modifies the density and temperature structure of gaseous halos, the formation of 
multiphase gas, and the dynamics of inflows and outflows. These effects produce a 
CGM that is smoother, more coherent, and more extended around massive galaxies, 
while suppressing cold gas around low-mass galaxies. The present Companion 
establishes the theoretical foundation for future studies of galaxy evolution, 
chemical enrichment, and the interaction between galaxies and the cosmic web in 
RDCC.

% ========================
% References
% ========================

\begin{thebibliography}{10}

\bibitem{Flagship} Lehmann, M. (2026). Relaxation-Driven Cyclic Cosmology (RDCC): A Minimal CPT-Symmetric Framework with a Single Parameter. Flagship Paper. Zenodo: \url{https://zenodo.org/records/19744958} (v43.0 or later).

\bibitem{Zenodo} The complete RDCC companion series (I--LVII) is available at the same Zenodo record above.

% Thematisch relevante Companions
\bibitem{CompXVI} Companion XVI: Boltzmann Code and Modified Hierarchy.
\bibitem{CompXXXI} Companion XXXI: RDCC Halo Model and Non-Linear Structure.
\bibitem{CompXXXII} Companion XXXII--XXXIX: Galaxy--Halo Connection, Cosmic Web, Lensing, RSD, ISW, tSZ, Rotation Curves, CGM.

% Wichtige externe Referenzen
\bibitem{Planck2020} Planck Collaboration (2020), Astron. Astrophys. 641, A6.
\bibitem{DESI2024} DESI Collaboration (2024/2025), arXiv: [aktuelle $f\sigma_8$/BAO-Paper].
\bibitem{JWST2024} Maiolino et al. (2024), Nature (JWST high-$z$ galaxies/quasars).
\bibitem{Kaiser1987} Kaiser, N. (1987), Mon. Not. R. Astron. Soc. 227, 1 (RSD).
\bibitem{Bartelmann2001} Bartelmann, M. \& Schneider, P. (2001), Phys. Rep. 340, 291 (Weak Lensing Review).
\bibitem{HuOkamoto2002} Hu, W. \& Okamoto, T. (2002), Astrophys. J. 574, 566 (CMB Lensing).
\bibitem{LewisChallinor2006} Lewis, A. \& Challinor, A. (2006), Phys. Rep. 429, 1 (CMB Lensing Review).
\bibitem{NFW1997} Navarro, J. F., Frenk, C. S. \& White, S. D. M. (1997), Astrophys. J. 490, 493 (Halo Profiles).

\end{thebibliography}

\end{document}
